% French PDF pp. 37--42.
\nde{Point (i) of Lemma 4.4.5 has been corrected, and the proofs of
the three points have been given in detail.}
Let \(R_\varnothing\) be the presheaf which associates the empty
set to every \(S\in\Ob\mathcal C\); it is an initial object of
\(\widehat{\mathcal C}\).

\begin{lemma}{4.4.5}\label{IV.4.4.5}
(i) Suppose that \(R_\varnothing\notin J(S)\), for every
\(S\in\Ob\mathcal C\). If \((X_i)\) is a family of separated
presheaves, the direct-sum presheaf \(\coprod_i X_i\) is separated.
\nde{In general, the direct sum of two sheaves \(F,G\) is not a
sheaf. Indeed, let \(S_1,S_2\in\Ob\mathcal C\); one assumes that
the direct sum \(S=S_1\amalg S_2\) exists in \(\mathcal C\) and
that the fibered product \(S_1\times_S S_2\) is an initial object
\(\varnothing\) of \(\mathcal C\) (cf. I, 1.8).
Let \(R\) be the sieve of \(S\) with basis \(\{S_1,S_2\}\);
then \((F\amalg G)(R)\) is the disjoint union of
\(F(S)\amalg G(S)\) and of \(F(S_i)\times G(S_j)\) for \(i\ne j\),
hence \(F\amalg G\) is not a sheaf in general.
On the other hand, if \(\mathcal C\) is the category with a single
object \(S\) and \(\mathrm{id}_S\) as its only morphism, endowed
with the topology defined by \(J(S)=\{R_\varnothing,S\}\),
then the only separated presheaves are \(R_\varnothing\) and
\(X=h_S\) (which is a sheaf), and \(X\amalg X\) is not separated.}

(ii) Consider an equivalence relation in the category of presheaves:
\[
\xymatrix{X\ar@<.5ex>[r]^u\ar@<-.5ex>[r]_v & Y }
\]
and let \(w:Y\to Z\) be the quotient.
If \(X\) is a sheaf and \(Y\) is separated, then \(Z\) is separated.

(iii) Consider an amalgamated sum in the category of presheaves,
where \(u,u'\) are monomorphisms:
\[
\xymatrix{
 & Y\ar[dr] & \\
X\ar[ur]^u\ar[dr]_{u'} & & V\\
 & Y'\ar[ur] & }.
\]
If \(Y,Y'\) are separated and \(X\) is a sheaf, then \(V\) is separated.
\end{lemma}
\begin{proof}
(i) Set \(X=\coprod_i X_i\). Let \(S\in\Ob\mathcal C\) and
\(\tau:R\hookrightarrow S\) be a refinement of \(S\), and let
\(x_1,x_2\) be two elements of \(X(S)\) such that
\(x_1\circ\tau=x_2\circ\tau\); there exist indices \(i,j\) such that
\(x_1\in X_i(S)\), \(x_2\in X_j(S)\).
Since \(R\ne R_\varnothing\), there exists a morphism
\(\phi:T\to R\), with \(T\in\Ob\mathcal C\).
Then \(x_1\circ\tau\circ\phi=x_2\circ\tau\circ\phi\), and since
\(X(T)\) is the disjoint union of the \(X_k(T)\), this implies \(i=j\).
Then, since \(X_i\) is separated and is a subobject of \(X\), the
map \(X_i(S)\to X_i(R)\to X(R)\) is injective, hence \(x_1=x_2\).
This proves that \(X\) is separated.

Let us prove (iii). Consider the morphisms \(i:Y\to V\) and
\(j:Y'\to V\), and let \(K\) be the kernel of:
\[
\xymatrix{Y\times Y'\ar@<.5ex>[r]^p\ar@<-.5ex>[r]_q & V },
\]
where \(p=i\circ\mathrm{pr}_1\), \(q=j\circ\mathrm{pr}_2\).
Let \(S\in\Ob(\mathcal C)\).
Since \(u,u'\) are monomorphisms, one may identify \(X(S)\) with
its image in \(Y(S)\) (resp. \(Y'(S)\)); denote by \(Z(S)\)
(resp. \(Z'(S)\)) the complement.
Then \(V(S)\) is identified with the disjoint union of
\(Z(S),Z'(S),X(S)\), and one easily sees that the maps
\(i(S):Y(S)\to V(S)\), \(j(S):Y'(S)\to V(S)\) are injective,
and that the map
\[
(u\times u')(S):X(S)\longrightarrow K(S)
\]
is bijective. Consequently \(i,j\) are monomorphisms, and
\(u\times u':X\to K\) is an isomorphism.

Set \(U=Y\amalg Y'\), and let \(\tau:R\hookrightarrow S\) be
a refinement of \(S\); one has a commutative diagram:
\[
\xymatrix{U(S)\ar[r]^{U(\tau)}\ar[d] & U(R)\ar[d]\\
V(S)\ar[r]_{V(\tau)} & V(R) }.
\]
Let \(v_1,v_2\in V(S)\) have the same image in \(V(R)\).
By definition of \(V\), \(v_1,v_2\) lift to elements \(y_1,y_2\)
of \(U(S)\); denote by \(z_1,z_2\) their images in \(U(R)\).
Then \(z_1,z_2\) have the same image in \(V(R)\).

Since \(i:Y\to V\), \(j:Y'\to V\) are monomorphisms, the maps
\(Y(R)\to V(R)\), \(Y'(R)\to V(R)\) are injective.
Thus, since \(Y,Y'\) are separated, if \(y_1,y_2\) both belong to
\(Y(S)\) or to \(Y'(S)\), then \(y_1=y_2\).
Otherwise one may assume that \(y_1\in Y(S)\), \(y_2\in Y'(S)\),
hence \(z_1\in Y(R)\), \(z_2\in Y'(R)\).
But then, since \(i\circ z_1=j\circ z_2\), the morphism
\(z_1\boxtimes z_2:R\to Y\times Y'\) factors through \(K=X\).
Since moreover \(X(R)=X(S)\) (because \(X\) is a sheaf), there
exists \(x\in X(S)\) such that
\(u(x)\circ\tau=z_1=y_1\circ\tau\) and
\(u'(x)\circ\tau=z_2=y_2\circ\tau\), whence, since \(Y,Y'\)
are separated, \(u(x)=y_1\), \(u'(x)=y_2\), and thus \(v_1=v_2\).
This proves that \(V\) is separated.

Let us prove (ii). First remark that the morphism of presheaves
\[
X\xrightarrow{u\boxtimes v}K,
\]
where \(K\) denotes the kernel of the pair of morphisms
\(w\circ\mathrm{pr}_i:Y\times Y\to Z\), is an isomorphism.
Indeed, for every \(T\in\Ob\mathcal C\), \(X(T)\) is an
equivalence relation in \(Y(T)\), so that the diagram
\[
\xymatrix{
X(T)\ar[r]^{u\boxtimes v} & Y(T)\times Y(T)
\ar@<.5ex>[r]^{w\circ\mathrm{pr}_1}
\ar@<-.5ex>[r]_{w\circ\mathrm{pr}_2} & Z(T) }
\]
is exact in the category of sets.

Now let \(S\in\Ob\mathcal C\) and \(g_1,g_2:S\to Z\) be two
morphisms which agree on a refinement \(\tau:R\hookrightarrow S\)
of \(S\). Since \(S\in\Ob\mathcal C\), one has
\(Z(S)=Y(S)/X(S)\), by construction of \(Z\), and hence there
exist morphisms \(f_1,f_2:S\to Y\) such that \(wf_i=g_i\).
Then \(wf_1\tau=wf_2\tau\), hence, by the preceding argument,
there exists a morphism \(\phi:R\to X\) such that
\(u\phi=f_1\tau\), \(v\phi=f_2\tau\).
Since \(X\) is a sheaf, there exists \(\psi:S\to X\) such that
\(\phi=\psi\tau\), and one therefore has in \(Y(R)\) the equalities:
\[
u\psi\tau=f_1\tau,\qquad v\psi\tau=f_2\tau.
\]
Since \(Y(S)\to Y(R)\) is injective (\(Y\) being separated),
this implies \(u\psi=f_1\), \(v\psi=f_2\), whence \(g_1=g_2\).
This proves that \(Z\) is separated.
\end{proof}

\begin{lemma}{4.4.6}\label{IV.4.4.6}
\nde{The statement of the lemma and its proof have been given in detail.}
(i) Let \(\{F_i\xrightarrow{f_i}F\}\) be a family of morphisms
of presheaves, and let \(G\subset F\) be the image presheaf.
Then the following diagram in \(\widehat{\mathcal C}\) is exact:
\[
\xymatrix{
(\coprod_i F_i)\times_G(\coprod_j F_j)
\ar@<.5ex>[r]^{\mathrm{pr}_1}\ar@<-.5ex>[r]_{\mathrm{pr}_2}
 & \coprod_i F_i\ar[r] & G }
\]
that is, for every presheaf \(H\), the following diagram of sets is exact:
\[
\xymatrix{
\Hom(G,H)\ar[r] & \prod_i\Hom(F_i,H)
\ar@<.5ex>[r]^p\ar@<-.5ex>[r]_q
 & \prod_{i,j}\Hom(F_i\times_G F_j,H) }.\tag{*}
\]

(ii) Every covering family of morphisms of sheaves is universal
effective epimorphic.
\end{lemma}
\begin{proof}
(i) Let \(H\) be a presheaf.
The map \(f^*\) which associates to a morphism \(\phi:G\to H\)
the family of morphisms \(\phi\circ f_i:F_i\to H\) is injective,
for, for every \(S\in\Ob\mathcal C\), \(\phi(S)\) is determined
by the \((\phi\circ f_i)(S)\), since the family
\(f_i(S):F_i(S)\to G(S)\) is surjective.
It is clear that the image of \(f^*\) is contained in \(\Ker(p,q)\).
Conversely, let \(\phi_i:F_i\to H\) be a family of morphisms
such that, for every \(i,j\), the diagram below is commutative:
\[
\xymatrix{F_i\times_G F_j\ar[r]\ar[d] & F_j\ar[d]^{\phi_j}\\
F_i\ar[r]_{\phi_i} & H }.
\]
Then, for every \(S\), the map \(\coprod_i F_i(S)\to H(S)\)
factors uniquely through a map \(\phi(S):G(S)\to H(S)\),
and this defines a morphism \(\phi:G\to H\) such that
\(f^*(\phi)=(\phi_i)\). This proves the exactness of sequence
(*), and point (i).

Let us prove (ii). Since the notion of covering family is stable
under base extension, it suffices to show that every covering
family is effective epimorphic. Thus let \(\{F_i\to F\}\) be a
covering family of morphisms of sheaves, and let \(G\) be the
image presheaf of this family. Since the family is covering,
the monomorphism \(G\hookrightarrow F\) is too, hence, by
4.3.12, one has \(a(G)=F\).

On the other hand, since \(G\hookrightarrow F\) is a monomorphism,
the fibered products \(F_i\times_G F_j\), \(F_i\times_F F_j\)
are the same. Thus, by (i), the following diagram of sets is exact,
for every presheaf \(H\):
\[
\xymatrix{
\Hom(G,H)\ar[r] & \prod_i\Hom(F_i,H)
\ar@<.5ex>[r]^p\ar@<-.5ex>[r]_q
 & \prod_{i,j}\Hom(F_i\times_F F_j,H) }.\tag{**}
\]
If \(H\) is moreover a sheaf, one has
\[
\Hom(G,H)=\Hom(a(G),H)=\Hom(F,H),
\]
and then (**) shows that \(\{F_i\to F\}\) is indeed an effective
epimorphic family in the category \(\widetilde{\mathcal C}\)
of sheaves.
\end{proof}

\begin{lemma}{4.4.7}\label{IV.4.4.7}
Every family of morphisms of sheaves \(\{F_i\to F\}\) factors
into a covering family \(\{F_i\to G\}\) and a monomorphism \(G\to F\).
\end{lemma}
\begin{proof}
Indeed, it suffices to take for \(G\) the image sheaf of the given family.
\end{proof}

\begin{proof}[Proof of Proposition 4.4.3]
One has seen in 4.4.6 that (iii) implies (ii), and evidently
(ii) implies (i). Finally let \(\{F_i\to F\}\) be an epimorphic
family; by Lemma 4.4.7, it factors into a covering family followed
by a monomorphism. But the latter, being dominated
\nde{We recall (cf. N.D.E. (17)) that a morphism \(g:G\to H\)
is said to be dominated by a family of morphisms \(F_i\to H\)
if this family factors through \(g\).}
by an epimorphic family, is an epimorphism, hence an isomorphism by 4.4.4.
\end{proof}

\begin{remark}{4.4.8}\label{IV.4.4.8}
\nde{To agree with later references, the numbering of the original,
which had two nos.\ 4.4.7, has been corrected.}
Since the image presheaf of the family \(\mathcal F\) is separated,
the construction of the associated sheaf shows that the conditions
of Proposition 4.4.3 are also equivalent to the following:

(v) For every \(S\in\Ob\mathcal C\), every \(f\in F(S)\) is
locally in the image of \(\mathcal F\), that is:

(vi) For every \(S\in\Ob\mathcal C\) and every \(f\in F(S)\),
the set of \(S'\to S\) such that the image of \(f\) in \(F(S')\)
is in the image of one of the \(F_i(S')\) is a refinement of \(S\).

(vii) (If the topology is defined by a pretopology \(R\).)
For every \(S\in\Ob\mathcal C\) and every \(f\in F(S)\), there
exists a family \(\{S_j\to S\}\in R(S)\) such that for every \(j\)
the image \(f_j\) of \(f\) in \(F(S_j)\) is in the image of one
of the \(F_i(S_j)\).
\end{remark}

\begin{remark}{4.4.8.bis}\label{IV.4.4.8.bis}
If the sheaf \(F\) is representable, the preceding conditions
are also equivalent to:

(viii) The set of \(T\to F\) (\(T\in\Ob\mathcal C\)) such that
there exist an \(i\) and a commutative diagram
\[
\xymatrix{F_i\ar[r] & F\\ T\ar[u]\ar[ur] & }
\]
is a refinement of \(F\).

Indeed, if (viii) is satisfied, the image presheaf of the \(F_i\)
is dominated by\nde{\og majorizes\fg{} has been replaced by
\og is dominated by\fg, cf. N.D.E. (17).} a refinement of \(F\),
which implies that the family is covering.
Conversely, one applies (vi) to \(\mathrm{id}_F\in F(F)\).

This condition is expressed in figurative language as follows:
locally on \(F\), there exists an \(i\) such that \(F_i\to F\)
possesses a section. In particular a morphism \(G\to F\), where
\(G\) is a sheaf and \(F\) a representable sheaf, will be covering
if and only if it possesses a section locally (on \(F\)).
\end{remark}

\nde{Lemmas 4.4.8.1 and 4.4.8.2 have been added.}
The following lemmas will be useful in VIB, 8.1 and 8.2.

\begin{lemma}{4.4.8.1}\label{IV.4.4.8.1}
Let \(Q\subset P\) be presheaves of groups on \(\mathcal C\),
with \(Q\) normal in \(P\). Suppose that \(P\) is separated
and that \(Q\) is a sheaf. Then the presheaf of groups \(P/Q\)
is separated. Consequently one has
\[
a(P/Q)=L(P/Q)=\varinjlim_{R\in J(S)}(P/Q)(R).
\]
\end{lemma}
\begin{proof}
It suffices to show the first assertion, for the second follows
from it by 4.3.11 (iv) and (ii).
Let \(S\in\Ob\mathcal C\), let \(R\) be a sieve of \(S\), and
let \(y\in P(S)/Q(S)\) have the identity as its image in \((P/Q)(R)\).
One must show that \(y=1\). Now, by hypothesis, the morphism
\(P(S)\to P(R)\) (resp. \(Q(S)\to Q(R)\)) is injective
(resp. an isomorphism), and in the commutative diagram below
the upper row is exact:
\[
\xymatrix@C=1.5em{
1\ar[r] & Q(S)\ar[r]\ar[d]_{\simeq} & P(S)\ar[r]\ar@{^{(}->}[d]
 & P(S)/Q(S)\ar[r]\ar[d] & 1\\
1\ar[r] & Q(R)\ar[r] & P(R)\ar[r] & (P/Q)(R) & }.
\]
The result will follow if one shows that the lower row is exact.
Let \(f:R\to P\) have the identity as its image in \((P/Q)(R)\),
and let \(\phi:T\to R\), with \(T\in\Ob\mathcal C\).
Then \(f\circ\phi\) is the identity of
\((P/Q)(T)=P(T)/Q(T)\), i.e. \(f\circ\phi\in Q(T)\).
Thus \(\phi\mapsto f\circ\phi\) is a functorial map
\(R(T)\to Q(T)\), hence defines a morphism of functors
\(\pi:R\to Q\) such that \(\pi(\mathrm{id}_R)=f\), whence \(f\in Q(R)\).
This proves the exactness of the lower row, and the lemma is proved.
\end{proof}

\begin{lemma}{4.4.8.2}\label{IV.4.4.8.2}
Let \(H\subset G\) be presheaves of groups on \(\mathcal C\).

(i) If \(H\) is normal in \(G\), then \(L(H)\) is normal in \(L(G)\).

(ii) If \(H\) is central in \(G\), then \(L(H)\) is central in \(L(G)\).
\end{lemma}
\begin{proof}
Let \(S\in\Ob\mathcal C\); one must show (cf. I 2.3.6) that
\(L(H)(S)\) is normal (resp. central) in \(L(G)(S)\).
Let \(h\in L(H)(S)\), \(g\in L(G)(S)\); there exist a sieve \(R\)
of \(S\) and elements \(h'\in H(R)\), \(g'\in G(R)\) such that
\(h=z_R(h')\), \(g=z_R(g')\) (notation of 4.3.10).
Since \(z_R\) is a morphism of groups,
\(ghg^{-1}h^{-1}=z_R(g'h'g'^{-1}h'^{-1})\).
In case (i), one has \(g'h'g'^{-1}h'^{-1}\in H(R)\), whence
\(ghg^{-1}h^{-1}\in LH(S)\); in case (ii),
\(g'h'g'^{-1}h'^{-1}=1\), hence \(ghg^{-1}h^{-1}=1\).
\end{proof}

\begin{proposition}{4.4.9}\label{IV.4.4.9}
Every equivalence relation in \(\widetilde{\mathcal C}\) is universal
effective (3.3.3): let \(R\) be a \(\widetilde{\mathcal C}\)-equivalence
relation in the sheaf \(X\); then the sheaf associated to the
separated presheaf
\[
i(X)/i(R):S\longmapsto X(S)/R(S)
\]
is a universal effective quotient of \(X\) by \(R\).
\end{proposition}
\begin{proof}
Let \(X/R\) be the quotient sheaf of \(X\) by \(R\), which exists
by 4.4.1 (ii): \(X/R=a(i(X)/i(R))\).
We must show that \(X\to X/R\) is a universal effective
epimorphism, and that the morphism \(f:R\to X\times_{X/R}X\)
is an isomorphism. The first assertion has already been proved
(4.4.3). As for \(f\), it comes, by application of the functor \(a\),
from the morphism \(i(R)\to i(X)\times_{i(X/R)}i(X)\) or,
since \(i(X)/i(R)\) is separated (4.4.5 (ii)), so that
\(i(X)/i(R)\to i(X/R)\) is a monomorphism, from the canonical
morphism \(i(R)\to i(X)\times_{i(X)/i(R)}i(X)\).

One is therefore reduced to proving the same assertion in the
category of presheaves. But \(i(X)/i(R)\) is the presheaf
\(S\mapsto X(S)/R(S)\), and one is reduced to proving the
analogous assertion in the category of sets, where it is immediate.
\end{proof}

\begin{proposition}{4.4.10}\label{IV.4.4.10}
Under the conditions of 4.4.9, let \(Y\) be a subsheaf of \(X\).
Denote by \(R_Y\) the equivalence relation induced in \(Y\) by \(R\).
Then the canonical morphism (3.1.6)
\[
Y/R_Y\longrightarrow X/R
\]
is a monomorphism: it identifies \(Y/R_Y\) with a subsheaf of
\(X/R\), which is the image sheaf of the composite morphism
\[
Y\longrightarrow X\longrightarrow X/R.
\]
\end{proposition}
\begin{proof}
The morphism of presheaves
\[
i(Y)/i(R_Y)=i(Y)/i(R)_{i(Y)}\longrightarrow i(X)/i(R)
\]
is a monomorphism. Since the functor \(a\) is left exact
(4.3.16), it transforms monomorphisms into monomorphisms, and thus
\[
Y/R_Y\longrightarrow X/R
\]
is a monomorphism. The last assertion follows from the commutative diagram
\[
\xymatrix{Y\ar[r]\ar[d] & X\ar[d]\\
Y/R_Y\ar[r] & X/R },
\]
and the fact that \(Y\to Y/R_Y\) is covering.
\end{proof}
By virtue of this proposition, we shall always identify \(Y/R_Y\)
with a subsheaf of \(X/R\).
